\bookchapter{第五十三章：数据存储与网络请求}

\chapterintro{一个只能“看”的 App 是死的；能存数据、能连网络的 App 才是活的。这一章认识安卓数据层的“三件套”：DataStore（存小配置）、Room（存结构化数据）、Retrofit（和服务器对话），外加它们的“总调度”——协程。学完这章，你就能写会记住用户习惯、会实时更新内容的真 App 了。}

\booksection{一、存储方案怎么选：先看这张表}

\begin{center}
\begin{tabularx}{\linewidth}{|X|X|X|X|}
\hline
\textbf{方案} & \textbf{适合存什么} & \textbf{特点} & \textbf{本书哪里用} \tabularnewline
\hline
DataStore & 零散配置（主题、开关） & 键值对，简单安全 & 第 54 章记账本设置 \tabularnewline
Room & 结构化数据（账目、用户） & 真正的 SQLite 数据库封装，可查询 & 第 54 章账目记录 \tabularnewline
Retrofit & （不存，联网取）服务器数据 & HTTP 请求库的事实标准 & 第 55 章天气数据 \tabularnewline
\hline
\end{tabularx}
\end{center}

\begin{center}\small\color{gray}数据三件套的分工\end{center}

记一个口诀：\textbf{“小配置进 DataStore，成批数据进 Room，外面的数据走 Retrofit”}。旧教程里常见的 SharedPreferences 已被 DataStore 取代，SQLite 手写 SQL 也已被 Room 取代——直接学新的，别走老路。

\booksection{二、DataStore：存点小配置}

想记住“用户上次选的币种”“是否开启深色模式”这类零散值，用 DataStore（Preferences 版）：

\begin{lstlisting}[language=Kotlin]
val Context.dataStore by preferencesDataStore(name = "settings")

// 写：存一个值
suspend fun saveTheme(isDark: Boolean) {
    context.dataStore.edit { prefs ->
        prefs[THEME_KEY] = isDark
    }
}

// 读：取出来（Flow 会自动通知变化）
val themeFlow: Flow<Boolean> = context.dataStore.data
    .map { it[THEME_KEY] ?: false }
\end{lstlisting}

注意返回的是\textbf{Flow}——一个会“流动”的数据流。值一旦改变，所有观察它的地方自动收到新值。这个概念（响应式）是安卓现代的基石，第 54 章会实战。

\booksection{三、Room：给数据安个“真正的家”}

记账本有成百上千笔账，还要按分类、按月份查询统计——这需要数据库。Room 是 SQLite 的官方封装，把“建表、查数据”变得像写 Kotlin 一样自然。三样东西：

\begin{lstlisting}[language=Kotlin]
// 1. Entity：表结构（一张表 = 一个 data class）
@Entity(tableName = "records")
data class Record(
    @PrimaryKey(autoGenerate = true) val id: Long = 0,
    val amount: Double,
    val category: String,
    val note: String,
    val date: String
)

// 2. Dao：操作接口（增删改查）
@Dao
interface RecordDao {
    @Query("SELECT * FROM records ORDER BY date DESC")
    fun getAll(): Flow<List<Record>>

    @Insert
    suspend fun insert(record: Record)

    @Delete
    suspend fun delete(record: Record)

    @Query("SELECT SUM(amount) FROM records WHERE category = :c")
    suspend fun sumByCategory(c: String): Double
}

// 3. Database：总装
@Database(entities = [Record::class], version = 1)
abstract class AppDatabase : RoomDatabase() {
    abstract fun recordDao(): RecordDao
}
\end{lstlisting}

这套“三件套”模板可以原样搬进你的任何项目——换换字段名就行。查询返回 Flow，界面跟着数据自动刷新，妙不可言。

\booksection{四、Retrofit：和服务器说话}

天气、新闻、汇率……这些数据不在手机里，在别人的服务器上。Retrofit 让“向服务器要数据”变得像调函数一样简单。三步走：

\begin{lstlisting}[language=Kotlin]
// 1. 定义接口：把 HTTP 请求描述成函数
interface WeatherApi {
    @GET("v1/weather")
    suspend fun getWeather(@Query("city") city: String): WeatherResponse
}

// 2. 创建实例
val api = Retrofit.Builder()
    .baseUrl("https://api.example.com/")
    .addConverterFactory(GsonConverterFactory.create())
    .build()
    .create(WeatherApi::class.java)

// 3. 像调函数一样调用（在协程里）
val result = api.getWeather("北京")
\end{lstlisting}

\textbf{避坑提醒：两条网络铁律}\par

(1)\textbf{网络请求绝不能在主线程做}——会卡死界面甚至崩溃（NetworkOnMainThreadException）。所有网络调用必须放在协程（viewModelScope.launch）里。(2) 明文 HTTP 默认被禁止：安卓 9 起要求 HTTPS。测试时若要用 http:// 的本地地址，需要在 Manifest 里显式声明\texttt{usesCleartextTraffic}——但正式产品必须全走 HTTPS。

\booksection{五、JSON 解析：给网络数据“翻译”}

服务器返回的数据大多是 JSON——一串像“嵌套字典”的文本。Retrofit 需要一个“翻译官”把它变成 Kotlin 对象。2026 年的两大主流方案：

\begin{center}
\begin{tabularx}{\linewidth}{|X|X|X|}
\hline
\textbf{方案} & \textbf{特点} & \textbf{推荐场景} \tabularnewline
\hline
Gson / Moshi & 老牌反射方案，配置零成本 & 快速起步、字段不多的小项目 \tabularnewline
kotlinx.serialization & 官方编译期方案，更快更安全 & 新项目首选（2026 年主流） \tabularnewline
\hline
\end{tabularx}
\end{center}

\begin{center}\small\color{gray}JSON 解析两大方案\end{center}

\begin{lstlisting}[language=Kotlin]
// 给数据类加 @Serializable 注解
@Serializable
data class WeatherResponse(
    val weather: List<Weather>,
    val main: Main,
    val name: String
)

// Retrofit 挂上对应 Converter（以 Retrofit 2.11+ 官方支持为例）：
// .addConverterFactory(json.asConverterFactory("application/json".toMediaType()))
\end{lstlisting}

无论哪种方案，铁律都是：\textbf{先看真实返回的 JSON，再定义字段}。很多“解析失败”的 bug，都是因为字段名写错、嵌套层级猜错——把原始 JSON 打印出来对照，十秒定位。

\booksection{六、协程：安卓并发的“普通话”}

\booksection{五、协程：安卓并发的“普通话”}

上面反复出现的\texttt{suspend}和\texttt{launch}来自 Kotlin 协程（Coroutine）——安卓官方的并发方案。你可以粗暴地理解为：\textbf{让“要等一会儿”的操作（网络、数据库）不阻塞界面}。

\begin{lstlisting}[language=Kotlin]
// 在 ViewModel 里启动协程
viewModelScope.launch {
    val records = dao.getAll().first()  // suspend：会自动挂起等待
    // 更新界面状态……
}

// 关键词速记：
// suspend   —— "我要等待，别拦着别人"
// launch    —— "开工！"
// viewModelScope —— 跟着页面走的"工作间"，页面销毁自动收工
\end{lstlisting}

零基础先记住用法就好，背后的“挂起与恢复”原理可以以后再深究。日常开发你写的 90\% 都是这个模板。

\begin{tcbpractice}{来动手：做一个"每日一句"小应用}
找一个免费的名言 API（比如 hitokoto 之类），用 Retrofit 拉取“每日一句”，显示在屏幕中央，加个“换一句”按钮（协程里重新请求）。再把喜欢的句子用 Room 存进“收藏夹”。这个练习把 DataStore/Room/Retrofit/协程全用了一遍——做完它，第 7、8 章就是小菜一碟。

\end{tcbpractice}

\begin{tcbsummary}{本章小结}
\begin{itemize}
  \item 口诀：小配置进 DataStore，成批数据进 Room，外面的数据走 Retrofit
  \item DataStore 返回 Flow：值变了，观察者自动收到
  \item Room 三件套：@Entity（表）+ @Dao（操作）+ @Database（总装），模板可直接复用
  \item Retrofit 三步：定义接口 创建实例 协程里调用
  \item 两条铁律：网络不进主线程；正式产品必须 HTTPS
  \item 协程 = 不阻塞界面的等待；viewModelScope 是页面专属工作间
\end{itemize}

\end{tcbsummary}
